\section{Validation des données} % Analyse exploratoire des données (Exploratory Data Analysis, EDA)
\label{section:validation_donnee}
Les variables $NASC$, $FTLE$, $FOD$ et $R_1, \ldots, R_9$ sont des séries spatio-temporelles, acquises le long des trajets des campagnes océanographiques. Plusieurs de leurs propriétés peuvent compromettre l'apprentissage ou fausser l'évaluation du modèle : un échantillonnage déséquilibré ou informatif, des différences de distribution entre sous-ensembles de données, et une dépendance entre observations voisines. Cette section vise à caractériser ces propriétés indépendamment de tout modèle, afin (1) de justifier les choix de prétraitement, (2) de construire un schéma de validation croisée adapté et (3) de délimiter le domaine dans lequel les estimations du modèle pourront être interprétées.

\subsection{Biais d'échantillonnage}
\label{subsection:biais_echantillonage}
Un biais d'échantillonnage peut affecter l'apprentissage de plusieurs manières. Par exemple : (1) une densité d'échantillonnage déséquilibrée conduit à un sous-apprentissage des conditions peu échantillonnées, et donc à une dégradation des performances prédictives dans ces conditions ; (2) lorsque la présence d'une observation dépend d'autres variables, voire de la valeur de la cible elle-même, les données disponibles ne sont plus représentatives de la population totale, ce qui induit un biais de sélection. Il est donc nécessaire d'identifier ces biais afin d'interpréter correctement les estimations du modèle et d'en encadrer la portée.

Les données de $NASC$ sont des données \textit{in situ}, acquises lors des campagnes océanographiques OBSAUSTRAL/THEMISTO. % À compléter : origine des covariables (FTLE, FOD, ratios pigmentaires) et unité d'échantillonnage (ESU)
On s'intéressera au déséquilibre de l'échantillonnage selon deux partitions naturelles des données : (1) la localisation spatiale et (2) la campagne océanographique (année et mois). Pour chacune de ces partitions, on étudiera (1) la densité d'échantillonnage et (2) la distribution et la nature des valeurs manquantes.

\paragraph{Densité d'échantillonnage}
Dans un premier temps, on produira des cartes de densité d'échantillonnage sur des grilles spatiales de résolutions différentes ($20 \times 20$~km et $1500 \times 1000$~km) afin d'identifier les déséquilibres spatiaux à différentes échelles. Le déséquilibre temporel sera décrit par des histogrammes du nombre d'unités élémentaires d'échantillonnage (\textit{Elementary Sampling Units}, ESU) par campagne océanographique. Les partitions seront ensuite croisées (par exemple, proportion d'ESU par zone et par campagne) : une modalité d'une partition échantillonnée exclusivement dans une modalité d'une autre rendrait leurs effets indiscernables pour le modèle.

On évaluera ensuite le caractère informatif de la densité d'échantillonnage, c'est-à-dire l'association entre le $NASC$ et la densité locale d'échantillonnage, calculée dans les fenêtres spatiales et temporelles considérées. Cette association sera quantifiée par une corrélation de rang de Spearman et par la comparaison des distributions du $NASC$ entre zones de forte et de faible densité. Une association marquée indiquerait que l'échantillonnage porte de l'information sur la cible, et donc qu'il existe un risque que le modèle effectue un \textit{shortcut learning}.

\paragraph{Distribution et nature des valeurs manquantes}
On quantifiera d'abord le taux de valeurs manquantes par variable et par partition naturelle des données. On cherchera ensuite à déterminer si certaines variables manquent conjointement, cette structure sera représentée par un graphique UpSet.

\subsection{Autocorrélation spatio-temporelle}
\label{subsection:autocorrelation}
Les observations acquises le long d'un trajet sont d'autant plus semblables qu'elles sont proches dans l'espace et dans le temps. Cette dépendance a trois conséquences : (1) elle réduit le nombre effectif d'observations indépendantes, ce qui rend optimistes les tests statistiques classiques; (2) elle compromet l'évaluation du modèle : lors du tirage bootstrap des Forêts Aléatoires, comme lors d'une validation croisée aléatoire, une observation de test a presque toujours des voisines dans le jeu d'entraînement, ce qui conduit à surestimer les performances (erreur \textit{out-of-bag} optimiste); (3) Elle favorise la mémorisation de la localisation : lorsque les covariables varient lentement dans l'espace, leur combinaison identifie une zone, et le modèle peut restituer le $NASC$ observé dans cette zone plutôt qu'apprendre une relation entre variables (\textit{shortcut learning}).

L'autocorrélation spatiale sera caractérisée par des variogrammes empiriques, et l'autocorrélation temporelle par la fonction d'autocorrélation le long des trajets ; l'indice de Moran fournira un résumé global. Ces analyses seront conduites pour le $NASC$ et pour chacune des covariables, dont les échelles de structuration peuvent différer. On en déduira la portée, c'est-à-dire la distance et la durée au-delà desquelles deux observations ne sont plus corrélées. Cette portée servira à dimensionner les blocs de la validation croisée, dont la taille devra lui être au moins égale, éventuellement complétée par une zone tampon entre blocs d'entraînement et de test.

\subsection{Décalage de distribution (\textit{Distribution shift})}
\label{subsection:distribution_shift}
Un modèle entraîné sur une partie des données puis appliqué à une autre (extrapolation) suppose que les deux sont issues d'une même distribution. Lorsque ce n'est pas le cas (\textit{distribution shift}), les performances mesurées sur les données d'entraînement ne sont pas représentatives des performances en conditions d'application. On étudiera ce décalage entre les modalités des partitions naturelles correspondant à l'usage visé du modèle, à savoir les campagnes et les zones géographiques.

Le décalage sera d'abord examiné sur les covariables. Pour chaque variable, les distributions seront comparées entre modalités (densités superposées, test de Kolmogorov-Smirnov), en portant une attention particulière aux plages de valeurs couvertes : les Forêts Aléatoires n'extrapolant pas au-delà de la plage observée à l'entraînement, une modalité dont les covariables sortent de cette plage se situerait hors du domaine d'applicabilité du modèle. Le décalage multivarié sera ensuite évalué par validation adversariale : un classifieur est entraîné à distinguer deux modalités à partir des seules covariables, et son aire sous la courbe ROC (AUC) mesure leur séparabilité, une valeur proche de $0{,}5$ indiquant des distributions indiscernables. L'importance des variables de ce classifieur identifie celles qui portent le décalage.

Le décalage sera également examiné sur la variable cible, en comparant les distributions du $NASC$ entre modalités (\textit{label shift}). La stabilité de la relation entre covariables et $NASC$, qui ne peut être évaluée sans modèle, sera étudiée lors de la validation du modèle.

Compte tenu du nombre d'observations et de leur potentielle autocorrélation, les tests statistiques seront presque systématiquement significatifs ; l'interprétation reposera donc sur la taille des effets (AUC, recouvrement des distributions) plutôt que sur les $p$-valeurs. Les décalages identifiés serviront à définir les blocs de la validation croisée, de sorte que l'évaluation reproduise les conditions de transfert réalistes, ainsi qu'à délimiter le domaine d'applicabilité du modèle.